\documentclass{article}
\usepackage{../assignment_preamble}

\title{Weekly Problem 9}
\author{Ravi Kini}
\date{June 1, 2026}

\begin{document}

\maketitle

\problem
\subproblem{(a)}
\begin{equation*}
\begin{nd}[rules=connexive,justformat=just]
\hypo {1} {\lnot (A \Rightarrow B)}
\open
\hypo {2} {A}
\have {3} {\lnot B} \nconde{1,2}
\have {4} {\lnot A \lor \lnot B} \oi{3}
\close
\open
\hypo {5} {\lnot A}
\have {6} {\lnot A \lor \lnot B} \oi{5}
\close
\have {7} {\lnot A \lor \lnot B} \tnd{2-4,5-6}
\end{nd}
\end{equation*}
We then see that $\lnot (A \Rightarrow B) \vdash \lnot A \lor \lnot B$ is a valid argument in connexive logic.

\subproblem{(b)}
Let $A = i$ and $B = 1$. We see that $\lnot A = i$ and $\lnot B = 0$, so $\lnot A \lor \lnot B = i$. However, $A \Rightarrow B = 1$, and so $\lnot(A \Rightarrow B) = 0$. Consequently, the argument $\lnot A \lor \lnot B \vdash \lnot (A \Rightarrow B)$ is an invalid argument in connexive logic.
\end{document}